\section{Open Questions}

The following five questions are genuinely open after this replication --- each is
motivated by a specific limit of the 2021 paper or of this replication itself, and
each has a concrete probe.

\begin{enumerate}

\item \textbf{Interaction with modern QND / dispersive readout.}
How does correlated-response readout mitigation compare to (and integrate with)
QND-optimized readout schemes on current IBM (Fez, Kingston) and Quantinuum
H-series devices, where hardware-level crosstalk suppression is already stronger
than 2021 Melbourne?
\emph{Probe:} run DDOT on a 15-qubit sub-block of each device, compare
$R_\text{true}$ vs tensor-product, quantify the correlated-vs-TP gap. Predict a
smaller-but-nonzero gap. Free IBM Quantum Open $+$ Quantinuum trial credits.

\item \textbf{Mid-circuit-measurement integration.}
Can the correlated inversion be applied \emph{between} mid-circuit measurements
without breaking downstream conditional-gate logic? Does the cluster structure
of readout crosstalk change when the classical branch is short vs long?
\emph{Probe:} 4-qubit circuit with one mid-circuit measurement gating an $X$ on
a neighbor. Simulate under the paper's noise model, then run on IBM dynamic
circuits. Compare end-of-circuit correlated inversion vs a proposed
\emph{forward-only} per-branch correction using the estimated true bit.

\item \textbf{Scaling to $N \geq 100$ qubits with $k$-local clusters.}
Does the method scale to 100+ qubits with $k=3$ or $k=4$ clusters, and where
does the $O(k\cdot 2^k \log N)$ DDOT budget actually break down (statistical
vs classical-inversion vs storage)?
\emph{Probe:} on a current 127-qubit IBM device, sweep $k\in\{2,3,4\}$ and
report (a) mean cluster fidelity residual vs $k$, (b) wall-clock
characterization time, (c) classical inversion time. The $k$ at which
residual plateaus is the chip's effective locality.

\item \textbf{ML-based cross-talk characterization.}
Can a GNN over the coupling map learn $R$ from a much smaller DDOT budget by
exploiting spatial locality and cross-cool-down correlations?
\emph{Probe:} collect $\geq 20$ daily DDOT snapshots from IBM's calibration
API, train a GNN with edges $=$ coupling map, node features $=$ current
single-qubit $T_1$/$T_2$/$p_{01}$/$p_{10}$/history. Test-set metric: mean
$L_1$ error of predicted vs measured cluster response. Baseline: use the
previous cool-down's response.

\item \textbf{Transferability across cool-down cycles.}
Half-life of the correlated calibration before mitigation quality degrades
below the tensor-product baseline?
\emph{Probe:} longitudinal study --- DDOT-characterize once, benchmark hourly
for 24--48 h without re-characterization. Plot $|\Delta\langle H\rangle|$
under (a) correlated with stale calibration, (b) TP with current-day
single-qubit calibration. Crossover time $=$ effective calibration half-life.
Free IBM Quantum Open access sufficient.

\end{enumerate}
